\addtocontents{toc}{\protect\Needspace{3\baselineskip}}
\originalchapter{闭弦层级匹配、引力子与耦合}\label{ch:basic-closed}
对应 8.251 第 20 讲和第 21 讲前半. 开弦的左右行波由反射关联, 闭弦的两组振子独立, 但两组零模位置与总动量不重复. 全局周期性进一步限制了允许态.

\originalsection{两组独立振子与一个零模}
在非紧致目标空间取 $0\leq\sigma<2\pi$, 自然单位下定义
\[
 X^I=x_0^I+\alpha'p^I\tau+
 \ii\sqrt{\frac{\alpha'}2}\sum_{n\ne0}\frac1n
 [\alpha_n^I\ee^{-\ii n(\tau-\sigma)}
 +\bar\alpha_n^I\ee^{-\ii n(\tau+\sigma)}].
\]
本编按 HW8.1 与 HW10.9 的约定, 无横杠模式依 $\tau-\sigma$, 横杠模式依 $\tau+\sigma$; 有的来源交换二者名称, 此时所有相应差式符号也交换. $\int P^{\tau I}\dd\sigma=p^I$ 要求零模系数为 $\alpha'$, 比开弦小一倍. 对两边振子均有
\[
 [\alpha_m^I,\alpha_n^J]=m\delta^{IJ}\delta_{m+n,0},
 \quad[\bar\alpha_m^I,\bar\alpha_n^J]=m\delta^{IJ}\delta_{m+n,0},
 \quad[\alpha_m,\bar\alpha_n]=0,
\]
而 $\alpha_0^I=\bar\alpha_0^I=\sqrt{\alpha'/2}p^I$. 只有一个 $x_0^I$ 和一个 $p^I$, 不应另给左右各一套可独立观测的位置零模.

\originalsection{局部约束与全局单值}
光锥规范取 $X^+=\alpha'p^+\tau$. 两边约束分别给
\[
 \sqrt{2\alpha'}\alpha_n^-=\frac{2L_n}{p^+},
 \qquad\sqrt{2\alpha'}\bar\alpha_n^-=\frac{2\bar L_n}{p^+},
\]
其中 $L_n=\frac12\sum_k:\alpha_{n-k}^I\alpha_k^I:$, 横杠同理.
经典零模在左右必须相同, 故 $L_0=\bar L_0$. 在量子正规序且每边截距为 1 的临界玻色弦,
\[
 L_0=\frac{\alpha'}4p^Ip^I+N,\qquad
 \bar L_0=\frac{\alpha'}4p^Ip^I+\bar N,
 \qquad N=\bar N.
\]
也可从 $X'^-=\dot X^IX'^I/(\alpha'p^+)$ 积分一圈得到
\[
 L_0-\bar L_0=-\frac{p^+}{2\pi}\int_0^{2\pi}X'^-\dd\sigma=0.
\]
负号来自本章已声明的左右命名. 最后等号假定非紧致 $X^-$ 严格周期, 若紧致目标允许绕数, 要重新加入相应零模差. 这个积分说明局部方程如何留下一个全局限制, 不是可任意删除的一半闭弦态.

\originalsection{哈密顿量与质量层}
两边正规序条件相加给
\[
 H_\tau=\alpha'p^+p^-=L_0+\bar L_0-2,
 \qquad M^2=\frac2{\alpha'}(N+\bar N-2).
\]
对层级匹配态 $N=\bar N$, $M^2=4(N-1)/\alpha'$. 基态 $N=0$ 是质量平方 $-4/\alpha'$ 的闭弦快子. 第一个无质量层 $N=\bar N=1$ 是
$R_{IJ}\alpha_{-1}^I\bar\alpha_{-1}^J|p\rangle$, 有 $24^2=576$ 态. 未匹配的 $N=1,\bar N=0$ 不属于此闭弦的物理谱.

\originalsection{引力子、二形式与膨胀子}
将矩阵 $R_{IJ}$ 唯一分解为对称无迹、反对称和迹,
\[
 R_{IJ}=\left[R_{(IJ)}-\frac{\delta_{IJ}}{D-2}R_{KK}\right]
 +R_{[IJ]}+\frac{\delta_{IJ}}{D-2}R_{KK}.
\]
这三部分分别对应 $D$ 维无质量引力子、反对称二形式 $B_{\mu\nu}$ 与标量膨胀子 $\Phi$. 对 $D=26$, 极化数为 $299+276+1=576$. 对一般 $D$ 的引力子极化公式为 $D(D-3)/2$, 其来源是无质量小群 $SO(D-2)$ 的对称无迹二阶张量, 不是把一个未施约束的 $D\times D$ 对称矩阵直接当物理极化.

场论中的同一信息体现在 $h_{\mu\nu}\mapsto h_{\mu\nu}+\partial_\mu\xi_\nu+\partial_\nu\xi_\mu$ 的线性规范冗余. 自由谱只识别这些粒子表示; 要产生非线性爱因斯坦引力, 还需检验弦相互作用及低能有效作用量. 第二编提供相应全息和低能工具.

\originalsection{弦耦合与膨胀子背景}
无量纲弦耦合取 $g_s=\ee^{\Phi_0}$. 对定向闭弦世界面, 一个额外把手使微扰权重增加 $g_s^2$, 三闭弦树级相互作用相应正比 $g_s$. 开弦的盘过程给 $g_o^2\propto g_s$, 因而开弦三点耦合按 $\sqrt{g_s}$ 标度, 具体常数随顶点与外态归一化变化.

膨胀子是动态场, 所以 $g_s$ 是某背景的常数值, 可随背景变化, 并非理论中永远不变的外部数字. 自由谱不能单独确定它的期望值. 对十维超弦按量纲和树级标度,
\[
 G_{10}\propto g_s^2\alpha'^4,\qquad
 \ell_{P,10}\propto g_s^{1/4}\ell_s.
\]
后一式用 $[G_{10}]=L^8$. 在弱耦合时普朗克长度小于弦长, 所以“弦长就是普朗克长度”只有在特定耦合数量级下才成立. 对未紧致化的 26 维玻色弦, 同类量纲式应改为 $G_{26}\propto g_s^2\alpha'^{12}$, 不把十维公式套到所有弦理论.

\originalsection{自由层与多弦粒子数}
一个振子基态加激发描述一条弦的内部状态. 对每个满足匹配的内部态, 低能场论可再建立创建该粒子的 Fock 算符, 描述多条弦. “一条弦有很多振子”与“很多条弦”是不同的计数问题; 对谱生成函数重复乘一遍全体 Fock 因子会混淆二者. 以下超弦和膜习题均先数单弦内部态.
